\begin{exercisechunk}
\item \textbf{Symbol renaming and external grounding.}
\label[exercise]{ex:l14-symbol-renaming}
\exerciseprofile{2}{2}{2}{\exproof\quad\exinterpretation}{%
Separate invariance under symbol renaming from the external relation that
assigns meanings to symbols.}
This exercise gives a precise model of formal symbol manipulation. The model
isolates what equality, copying, and lookup preserve, and the final parts state
exactly what additional information a grounding relation supplies.

Let \(\cS\) be a finite set of symbols. A \emph{symbolic input} \(x\) is an
ordered finite collection of symbols together with zero or more finite tables
of ordered pairs from \(\cS\times\cS\). A deterministic \emph{symbolic
program} \(\cA\) has symbol registers and Boolean registers, contains no
hard-coded element of \(\cS\), and may perform only the following operations:
\begin{enumerate}[label=(\roman*),leftmargin=*]
\item copy a symbol from an input location or symbol register;
\item test whether two symbol registers contain the same symbol;
\item use a symbol as a key in an input table with distinct keys and copy the
corresponding value; and
\item branch according to a Boolean register.
\end{enumerate}
Assume that \(\cA\) terminates and returns one symbol in \(\cS\) on every
valid input.

Let \(\pi:\cS\to\cS\) be a bijection. Define \(\pi x\) by replacing every
symbol \(s\) occurring in \(x\), including every table key and value, by
\(\pi(s)\), while preserving positions and table order. For an output symbol
\(s\), the expression \(\pi(s)\) has its ordinary meaning: apply the map
\(\pi\) to \(s\). The program \(\cA\) itself is unchanged.
\begin{enumerate}[(a),beginpenalty=10000]
\item Prove that copying, equality testing, and table lookup commute with the
renaming \(\pi\). State the corresponding relation for each operation.
\item Prove by induction over the executed operations that
\[
 \cA(\pi x)=\pi\bigl(\cA(x)\bigr)
 \qquad\text{for every valid symbolic input }x.
\]
\item Let \(\cS=\{s_0,s_1\}\) and let
\(\cO=\{o_0,o_1\}\) be a set of external objects. A \emph{grounding} is a
bijection \(g:\cS\to\cO\). Consider
\[
 g_0(s_i)=o_i,
 \qquad
 g_1(s_i)=o_{1-i},
 \qquad i\in\{0,1\}.
\]
Fix one symbolic input \(x\), and suppose \(\cA\) receives neither \(g_0\)
nor \(g_1\). Prove that \(\cA(x)\) cannot name the target object \(o_0\)
under both groundings; that is, at least one of
\[
 g_0(\cA(x))=o_0,
 \qquad
 g_1(\cA(x))=o_0
\]
must fail.
\item Explain how parts~(b) and (c) differ. Your explanation must identify
what the renaming result proves, why the impossibility in part~(c) requires
the grounding to be absent from the input, and what kind of additional input
could break the ambiguity.
\end{enumerate}
\end{exercisechunk}
